\chapter*{RESUMEN}
\addcontentsline{toc}{chapter}{Resumen}
\markboth{RESUMEN}{RESUMEN}

El constante crecimiento del Internet de las Cosas (IoT) en las empresas enfrenta graves vulnerabilidades, ya que los Sistemas de Detección de Intrusiones (IDS) tradicionales tienen muy poca sensibilidad para detectar ciberataques que son infrecuentes pero sumamente peligrosos, como aquellos basados en web o de fuerza bruta. Para solucionar esto, esta investigación desarrolló y evaluó un modelo híbrido de clasificación apoyado en una arquitectura Stacking de dos niveles, combinando los algoritmos Random Forest, XGBoost GPU y LightGBM. Todo esto fue reforzado con la Entropía de Shannon para medir la incertidumbre y la técnica Borderline-SMOTE1 para equilibrar los datos. Este diseño cuasi-experimental se aplicó sobre el dataset CICIoT2023, analizando más de 2.5 millones de flujos de red totalmente limpios, destinando el 80\% de la información para el entrenamiento y aislando el 20\% exclusivamente para las pruebas. Los resultados demostraron el éxito de la propuesta al lograr un F1-Macro del 70.18\% y una exactitud del 84.85\%, superando con claridad a todos los modelos individuales evaluados. Además, registró un desempeño asombroso con una latencia de apenas $8.65\,\mu s$ por flujo y un consumo mínimo de $257.89\text{ MB}$ de memoria RAM, garantizando que puede funcionar sin problemas en pasarelas de borde. Finalmente, tras confirmar su validez matemática mediante la Prueba de Wilcoxon ($p = 0.003744 < 0.05$), concluimos con certeza que mezclar estos clasificadores resulta ser una solución sumamente eficaz, rápida y liviana para proteger las redes IoT.

\vspace{0.5cm}
\noindent \textbf{Palabras clave:} Ataques minoritarios, Borderline-SMOTE, CICIoT2023, Ciberseguridad, Entropía de Shannon, Internet de las Cosas, Stacking híbrido.

\clearpage
\chapter*{ABSTRACT}
\addcontentsline{toc}{chapter}{Abstract}
\markboth{ABSTRACT}{ABSTRACT}

The constant growth of the Internet of Things (IoT) in businesses faces serious vulnerabilities, as traditional Intrusion Detection Systems (IDS) have very low sensitivity when it comes to detecting cyberattacks that are infrequent but highly dangerous, such as web-based or brute force attacks. To address this issue, this research developed and evaluated a hybrid classification model based on a two-level Stacking architecture, combining the Random Forest, XGBoost GPU, and LightGBM algorithms. This approach was reinforced with Shannon Entropy to measure uncertainty and the Borderline-SMOTE1 technique to balance the data. This quasi-experimental design was applied to the CICIoT2023 dataset, analyzing over 2.5 million fully cleansed network flows, strategically allocating 80\% of the information for training and isolating the remaining 20\% exclusively for testing. The results demonstrated the success of the proposal by achieving an F1-Macro of 70.18\% and an accuracy of 84.85\%, clearly outperforming all the individual models evaluated. Furthermore, it recorded an astonishing computational performance with a latency of just $8.65\,\mu s$ per flow and a minimal consumption of 257.89 MB of RAM, ensuring it can operate smoothly on edge gateways. Finally, after confirming its mathematical validity through the Wilcoxon Test ($p = 0.003744 < 0.05$), we can conclude with absolute certainty that combining these classifiers provides a highly effective, fast, and lightweight solution for protecting IoT networks.

\vspace{0.5cm}
\noindent \textbf{Keywords:} Borderline-SMOTE, CICIoT2023, Cybersecurity, Hybrid Stacking, Internet of Things, Minority Attacks, Shannon Entropy.
